\element{loiHess}{
  \begin{question}{loiHess}\bareme{1}
    Calculer l'enthalpie de réaction de la combustion du méthane par du dioxygène pour former du dioxyde de carbone et de l'eau.

    \begin{tabular}{|l|c|c|c|}
      \hline
      Espèces chimiques                            & \ce{CH4(g)}    & \ce{CO2(g)}     & \ce{H2O(g)}  \\
      \hline
      $\Delta_fH^\circ_m$ (\unit{kJ.mol^{-1}})     & \num{-74.9}    & \num{-393.52}   & \num{-241.8} \\
      \hline
    \end{tabular}
    
    \begin{multicols}{4}
      \begin{reponses}
        \bonne{\SI{-802.2}{kJ/mol}}
        \mauvaise{\SI{802.2}{kJ/mol}}
        \mauvaise{\SI{-560.4}{kJ/mol}}
        \mauvaise{\SI{560.4}{kJ/mol}}
      \end{reponses}
    \end{multicols}
  \end{question}
}

\element{exothermique}{
  \begin{questionmult}{exothermique}\bareme{haut=2,p=0}
    Une réaction chimique d'enthalpie de réaction \SI{230}{kJ/mol} et d'entropie de réaction \SI{-397}{J.K^{-1}.mol^{-1}}
    % \begin{multicols}{4}
      \begin{reponses}
        \bonne{Consomme davantage de gaz qu'elle n'en produit}
        \mauvaise{Est exothermique}
        \mauvaise{Libère de la chaleur lorsqu'elle est faite de façon isotherme}
        \bonne{Fait baisser la température lorsqu'elle est faite de façon adiabatique}
      \end{reponses}
    % \end{multicols}
  \end{questionmult}
}
